
A commonly stated explanation for RLEF's advantage over SFT on hard competitive programming benchmarks is that the training dataset lacks correct solutions at competition difficulty, leaving SFT with near-zero gradient signal. This explanation predicts that APPS, the most widely used algorithmic training dataset, should contain sparse reference solution coverage at competition level.

Empirical measurement contradicts this prediction. The APPS dataset contains reference solutions passing all included test cases for 85.32\% of competition-level problems (308/361). Yet a DeepSeek-Coder-7B model fine-tuned with SFT on those solutions achieves 0.0 pass@1 on LiveCodeBench-Hard. The dataset is not sparse; the model cannot generalize the solutions present in training to held-out hard evaluation problems. This observation --- a generalization void rather than a dataset void --- is the empirical starting point of this paper.

This reframing has consequences for how RLEF's advantage should be understood. If SFT's failure at hard benchmarks is a generalization problem rather than a data coverage problem, then RLEF's benefit may derive from training on the model's own generated outputs at its actual capability frontier, rather than from overcoming missing training data. This hypothesis further predicts that the specific reward formulation --- which governs the granularity of the execution signal --- should matter less than the existence of execution feedback itself.

The reward formulation prediction is testable and is tested here. Fraction-of-tests and binary RLEF rewards produce statistically equivalent performance in the controlled experiment reported in this paper (p=0.552). The execution infrastructure matters; the reward function granularity does not, at the training scale examined.

Despite the practical importance of RLEF for code generation, no prior work provides a controlled, reproducible, open-source comparison of RLEF against SFT across the full benchmark difficulty spectrum from HumanEval to LiveCodeBench-Hard. Existing work either uses non-public models \citep{Gehring2024RLEF} or evaluates only on easy-to-medium difficulty benchmarks (Le et al., 2022; Majeed et al., 2023; Liu et al., 2023). This reproducibility gap prevents verification of difficulty-scaling claims.

The present study fills this gap using publicly available components throughout: DeepSeek-Coder-7B as the base model, APPS for training, and bigcode-evaluation-harness for evaluation across five difficulty levels. All variables except training method (SFT versus RLEF) are held fixed.

\textbf{Contributions:}

\begin{enumerate}
\item \textbf{Controlled reproducible pipeline.} An open-source RLEF versus SFT comparison framework using DeepSeek-Coder-7B, APPS, bigcode-harness, and a custom SimpleGRPOTrainer (compatible with PyTorch 2.5.1 / TRL 1.x), validated through a multi-hypothesis experimental protocol spanning five sub-experiments.
\end{enumerate}

\begin{enumerate}
\item \textbf{Difficulty-scaling evidence.} A statistically significant positive-ordered trend of RLEF's advantage over SFT across five benchmark difficulty levels (JT z=+56.10, $p\approx 0)$, with the largest gains at LiveCodeBench-Hard ($\Delta =+0.18$, smoke-scale proxy) and SFT achieving 0.0 pass@1 at that level (directly measured).
\end{enumerate}

\begin{enumerate}
\item \textbf{Null result on reward formulation.} A controlled negative result demonstrating that fraction-of-tests and binary RLEF rewards converge to statistically equivalent performance at APPS training scale (p=0.552), consistent with and independently replicating two findings in the literature on a different model and dataset combination.
\end{enumerate}

\begin{enumerate}
\item \textbf{Generalization void characterization.} Empirical demonstration that 85.32\% APPS competition-level reference solution coverage coexists with 0.0 SFT pass@1 on LiveCodeBench-Hard, decoupling dataset coverage from model generalization ability and motivating a revised mechanistic framing of RLEF's hard-difficulty advantage.
\end{enumerate}

